% ---------------------------------------------------------------------
%  CHƯƠNG 2 — CƠ SỞ LÝ THUYẾT VÀ CÔNG NGHỆ
% ---------------------------------------------------------------------
\chapter{Cơ sở lý thuyết và công nghệ}
\label{ch:coso}

\begin{muctieu}
Chương này trình bày nền tảng lý thuyết mà hệ thống dựa vào: cách một khuôn
mặt được phát hiện và chuyển thành vector đặc trưng, cách so khớp danh tính
bằng khoảng cách cosine, nguyên lý chống giả mạo khuôn mặt, kỹ thuật tìm
kiếm trên không gian vector, cùng các công nghệ được chọn để hiện thực hoá.
\end{muctieu}

\section{Bài toán nhận diện khuôn mặt}

Cần phân biệt ba bài toán thường bị gộp làm một \cite{Jain2004Biometric, Wang2021Survey}:

\begin{dinhnghia}[Ba bài toán trên khuôn mặt]
\textbf{Phát hiện khuôn mặt} \en{(face detection)} trả lời câu hỏi trong ảnh
có khuôn mặt nào và nằm ở đâu. \textbf{Xác minh khuôn mặt} \en{(face
verification)} trả lời hai khuôn mặt có phải cùng một người hay không.
\textbf{Nhận diện khuôn mặt} \en{(face recognition hay identification)} trả
lời khuôn mặt này là ai trong số những người đã biết.
\end{dinhnghia}

Điểm danh trong lớp là bài toán nhận diện trên tập đóng có giới hạn: hệ thống
chỉ cần xác định khuôn mặt trong khung hình thuộc về sinh viên nào trong danh
sách lớp, chứ không phải toàn bộ cơ sở dữ liệu của trường. Việc thu hẹp
không gian tìm kiếm theo lớp học phần vừa làm tăng tốc độ, vừa giảm mạnh xác
suất nhận nhầm, và được hệ thống khai thác triệt để ở Mục~\ref{sec:sokhop}.

Quy trình nhận diện hiện đại gồm bốn bước nối tiếp nhau \cite{Wang2021Survey} như Hình~\ref{fig:pipeline}.

\begin{figure}[H]
\centering
\begin{tikzpicture}[node distance=1.15cm and 0.85cm, every node/.append style={align=center}]
  \node (anh)   [io]      {Khung hình\\ từ camera};
  \node (detect)[process, right=of anh]    {Phát hiện\\ SCRFD};
  \node (align) [process, right=of detect] {Căn chỉnh\\ theo 5 điểm mốc};
  \node (embed) [process, right=of align]  {Trích đặc trưng\\ ArcFace};
  \node (match) [decision, below=of align] {So khớp cosine};
  \node (out)   [startstop, right=of match, xshift=0.6cm] {Danh tính\\ và độ tin cậy};
  \draw [flowarr] (anh)    -- (detect);
  \draw [flowarr] (detect) -- (align);
  \draw [flowarr] (align)  -- (embed);
  \draw [flowarr] (embed)  |- (match);
  \draw [flowarr] (match)  -- (out);
\end{tikzpicture}
\caption{Bốn bước của quy trình nhận diện khuôn mặt trong FAMS.}
\label{fig:pipeline}
\end{figure}

\section{Phát hiện khuôn mặt bằng SCRFD}

Bước đầu tiên là tìm vị trí mọi khuôn mặt trong khung hình. Các phương pháp
cổ điển dựa trên đặc trưng Haar và bộ phân loại tầng \cite{Viola2001Rapid}
chạy nhanh nhưng suy giảm rõ rệt khi khuôn mặt nghiêng hoặc thiếu sáng, tức
đúng những điều kiện thường gặp trong giảng đường.

Hệ thống dùng SCRFD \cite{Guo2021SCRFD}, một bộ dò khuôn mặt một giai đoạn
được thiết kế để phân bổ chi phí tính toán hợp lý giữa các tầng đặc trưng,
nhờ đó đạt độ chính xác cao trên khuôn mặt nhỏ mà vẫn nhẹ. Với ảnh chụp cả
lớp, khuôn mặt của sinh viên ngồi cuối phòng chiếm rất ít điểm ảnh, nên khả
năng phát hiện khuôn mặt nhỏ là tiêu chí quyết định lựa chọn này. Ngoài khung
bao, SCRFD trả về năm điểm mốc gồm hai mắt, chóp mũi và hai khoé miệng, dùng
cho bước căn chỉnh tiếp theo.

Căn chỉnh \en{(alignment)} là phép biến đổi affine đưa năm điểm mốc về vị trí
chuẩn, cho ra một ảnh khuôn mặt kích thước cố định \cite{Wang2021Survey}. Bước này quan trọng vì mô
hình trích đặc trưng được huấn luyện trên các khuôn mặt đã căn chỉnh; bỏ qua
nó thì cùng một người ở hai góc chụp khác nhau sẽ cho hai vector lệch nhau
đáng kể.

\section{Biểu diễn khuôn mặt bằng ArcFace}
\label{sec:arcface}

Thay vì phân loại trực tiếp ảnh khuôn mặt thành các lớp là từng người, hướng
tiếp cận hiện đại học một \term{hàm nhúng} \en{(embedding)} ánh xạ ảnh khuôn
mặt vào không gian vector, sao cho hai ảnh của cùng một người nằm gần nhau và
hai người khác nhau nằm xa nhau. Ý tưởng này được phổ biến bởi FaceNet
\cite{Schroff2015FaceNet} và được ArcFace \cite{Deng2019ArcFace} cải tiến.

ArcFace bổ sung một biên góc cộng tính vào hàm mất mát softmax. Gọi
$\theta_{y_i}$ là góc giữa vector đặc trưng của mẫu thứ $i$ và vector trọng
số của lớp đúng $y_i$, $s$ là bán kính cầu và $m$ là biên góc, hàm mất mát
có dạng \cite{Deng2019ArcFace}:

\begin{equation}
\label{eq:arcface}
L = -\frac{1}{N}\sum_{i=1}^{N}
\log \frac{e^{s\cos(\theta_{y_i}+m)}}
{e^{s\cos(\theta_{y_i}+m)} + \sum_{j \neq y_i} e^{s\cos\theta_j}}
\end{equation}

Việc cộng thêm $m$ vào góc của lớp đúng buộc mô hình phải kéo mẫu lại gần
tâm lớp của nó một khoảng chặt hơn mức tối thiểu cần để phân loại đúng. Hệ
quả là các cụm đặc trưng của từng người tách bạch hơn trên mặt cầu đơn vị, và
đó chính là tính chất mà bài toán nhận diện cần.

Mô hình được dùng là bộ \code{buffalo\_l} của thư viện InsightFace
\cite{InsightFace2023}, sinh ra vector \FchieuVector{} chiều đã chuẩn hoá về
độ dài đơn vị.

\begin{ghinho}
Vector đặc trưng đã chuẩn hoá nên tích vô hướng của hai vector chính bằng
cosine của góc giữa chúng. Nhờ vậy phép so khớp danh tính chỉ còn là một
phép nhân vô hướng, rất rẻ về mặt tính toán.
\end{ghinho}

\subsection{So khớp bằng độ tương tự cosine}
\label{sec:sokhop}

Với hai vector đặc trưng $\mathbf{u}$ và $\mathbf{v}$, độ tương tự cosine và
khoảng cách cosine được định nghĩa \cite{Manning2008IR}:

\begin{equation}
\label{eq:cosine}
\mathrm{sim}(\mathbf{u}, \mathbf{v}) =
\frac{\mathbf{u} \cdot \mathbf{v}}{\lVert \mathbf{u}\rVert \, \lVert \mathbf{v}\rVert},
\qquad
d(\mathbf{u}, \mathbf{v}) = 1 - \mathrm{sim}(\mathbf{u}, \mathbf{v})
\end{equation}

Giá trị $\mathrm{sim}$ nằm trong đoạn $[-1, 1]$, càng gần $1$ thì hai khuôn
mặt càng giống nhau. Hệ thống coi là cùng một người khi độ tương tự vượt
ngưỡng $\tau = \FnguongNhanDien$. Ngưỡng này điều tiết trực tiếp cán cân giữa
hai loại lỗi \cite{Jain2004Biometric, ISO19795}:

\begin{itemize}[leftmargin=1.6em, itemsep=2pt]
  \item \textbf{Chấp nhận sai} \en{(false accept)}: nhận nhầm sinh viên này
        thành sinh viên khác, dẫn tới điểm danh cho người vắng mặt. Đây là
        lỗi nguy hiểm vì phá hỏng ý nghĩa của cả hệ thống.
  \item \textbf{Từ chối sai} \en{(false reject)}: không nhận ra sinh viên
        đang có mặt. Lỗi này gây phiền nhưng còn khắc phục được, vì giảng
        viên có thể xác nhận tay.
\end{itemize}

Vì hai loại lỗi có hậu quả không cân xứng, hệ thống ưu tiên hạ thấp tỉ lệ
chấp nhận sai, chấp nhận đánh đổi bằng một mức từ chối sai cao hơn.

\subsection{Vì sao không tự huấn luyện mô hình}
\label{sec:vosaokhonghuanluyen}

Một câu hỏi hợp lý là tại sao đồ án không huấn luyện một mạng phân loại trên
chính ảnh sinh viên của trường, mà dùng lại mô hình đã huấn luyện sẵn theo
hướng học chuyển giao \en{(transfer learning)} \cite{Pan2010Transfer}. Có ba lý do:

\begin{enumerate}[leftmargin=1.8em, itemsep=2pt]
  \item \textbf{Tập người dùng luôn thay đổi.} Mỗi kỳ có sinh viên mới nhập
        học và sinh viên tốt nghiệp. Mô hình phân loại buộc phải huấn luyện
        lại mỗi lần danh sách đổi, còn hàm nhúng chỉ cần thêm một vector.
  \item \textbf{Dữ liệu cho mỗi người quá ít.} Mỗi sinh viên chỉ đăng ký
        vài ảnh, không đủ để huấn luyện một lớp phân loại ổn định, trong khi
        ArcFace được huấn luyện trên hàng triệu ảnh của hàng chục nghìn danh
        tính \cite{Deng2019ArcFace} nên đã học được đặc trưng tổng quát.
  \item \textbf{Chi phí và rủi ro.} Huấn luyện lại đòi hỏi GPU và một tập
        dữ liệu khuôn mặt sinh viên đủ lớn, tức phải thu thập thêm dữ liệu
        sinh trắc, đi ngược nguyên tắc thu thập tối thiểu \cite{GDPR2016} đã nêu ở
        Chương~\ref{ch:modau}.
\end{enumerate}

\section{Chống giả mạo khuôn mặt}
\label{sec:chonggiamao}

Một hệ thống chỉ so khớp đặc trưng sẽ chấp nhận cả ảnh in lẫn ảnh hiển thị
trên màn hình điện thoại, vì những ảnh đó cho vector đặc trưng gần như giống
hệt người thật. Đây là kiểu tấn công trình chiếu \en{(presentation attack)} \cite{ISO30107, Ramachandra2017Survey},
và là lỗ hổng rõ ràng nhất của điểm danh bằng khuôn mặt.

Phát hiện sự sống \en{(liveness detection)} chia làm hai nhóm \cite{Ramachandra2017Survey}. Nhóm chủ động
yêu cầu người dùng thực hiện hành vi như chớp mắt hay quay đầu; nhóm bị động
chỉ phân tích chính khung hình tĩnh. Hệ thống chọn hướng bị động với mô hình
MiniFASNetV2 thuộc bộ Silent-Face-Anti-Spoofing \cite{Minivision2020Silent},
vì hướng chủ động không dùng được khi giảng viên quét cả lớp bằng một khung
hình.

MiniFASNetV2 khai thác những dấu vết mà một bề mặt phẳng để lại \cite{Ramachandra2017Survey}: phản xạ ánh
sáng khác da thật, hiện tượng răng cưa của lưới điểm ảnh khi chụp lại màn
hình, viền giấy hoặc viền thiết bị, và sự thiếu vắng kết cấu vi mô của da.
Mô hình nhận ảnh khuôn mặt đã mở rộng khung bao theo hệ số $2{,}7$ để lấy
thêm bối cảnh xung quanh, đưa về kích thước $80 \times 80$ điểm ảnh, rồi cho
ra điểm tin cậy cho ba lớp, trong đó lớp \term{mặt thật} phải vượt ngưỡng
$\FnguongThat$ thì khung hình mới được chấp nhận.

\begin{luuy}
Không có mô hình chống giả mạo nào đạt độ chính xác tuyệt đối. Vì vậy hệ
thống không coi đây là lớp phòng vệ duy nhất, mà kết hợp thêm xác nhận đa
khung hình, đối chiếu vị trí và quyền xác nhận cuối cùng của giảng viên, như
trình bày ở Mục~\ref{sec:baolop}.
\end{luuy}

\section{Tìm kiếm trên không gian vector}

Khi số vector đã đăng ký tăng lên, việc so sánh tuần tự với từng vector trở
nên tốn kém. Bài toán này gọi là tìm láng giềng gần nhất \en{(nearest
neighbour search)}. Với dữ liệu nhiều chiều, các cấu trúc chỉ mục cổ điển
không còn hiệu quả, nên thực tế dùng tìm kiếm gần đúng \cite{Indyk1998ANN}, đổi một phần rất nhỏ
độ chính xác lấy tốc độ cao hơn nhiều lần.

Hệ thống dùng chỉ mục HNSW \cite{Malkov2020HNSW}, một đồ thị thế giới nhỏ
phân tầng: tầng trên thưa để nhảy nhanh tới vùng cần tìm, tầng dưới dày để
tinh chỉnh kết quả. Điểm thuận lợi về mặt kỹ thuật là chỉ mục này có sẵn
trong pgvector \cite{Pgvector2024}, một phần mở rộng của PostgreSQL cho phép
lưu vector ngay trong bảng dữ liệu và viết truy vấn so khớp bằng SQL thông
thường. Nhờ đó hệ thống không cần thêm một cơ sở dữ liệu vector riêng, và
điều kiện lọc theo lớp học phần hay theo trạng thái duyệt được đẩy xuống cho
cơ sở dữ liệu xử lý cùng lúc với phép so khớp.

\section{Công nghệ sử dụng}

\begin{table}[H]
\centering
\caption{Các công nghệ chính và lý do lựa chọn.}
\label{tab:congnghe}
\small
\begin{tabularx}{\textwidth}{@{}>{\bfseries}l >{\raggedright\arraybackslash}p{3.1cm} X@{}}
\toprule
Thành phần & Công nghệ & Lý do chọn \\
\midrule
Dịch vụ & FastAPI, SQLAlchemy 2.0 & Kiểu dữ liệu tường minh, tự sinh tài liệu
  API, hệ sinh thái Python sẵn có thư viện thị giác máy tính \\
Cơ sở dữ liệu & PostgreSQL 16, pgvector \cite{Pgvector2024} & Lưu dữ liệu quan hệ và vector đặc
  trưng trong cùng một nơi, hỗ trợ chỉ mục HNSW \\
Suy luận mô hình & ONNX Runtime (CPU) & Chạy được trên máy chủ không GPU,
  không phụ thuộc khung học sâu lúc triển khai \\
Ứng dụng & Flutter & Một mã nguồn cho Web, Android và iOS, truy cập được
  camera và định vị \\
Triển khai & Docker Compose, nginx & Dựng lại toàn bộ hệ thống bằng một
  lệnh, tách biệt môi trường \\
Kiểm thử & pytest trên PostgreSQL thật & Kiểu dữ liệu vector không mô phỏng
  được bằng cơ sở dữ liệu nhúng \\
\bottomrule
\end{tabularx}
\end{table}

{\small
\begin{xltabular}{\textwidth}{@{}>{\bfseries\raggedright\arraybackslash\leavevmode}p{2.6cm} >{\raggedright\arraybackslash\leavevmode}p{5.6cm} >{\raggedright\arraybackslash\leavevmode}X@{}}
\caption{Giải thích chi tiết Bảng~\ref{tab:congnghe}: mỗi công nghệ nằm ở đâu trong mã nguồn.}
\label{tab:congnghe-gt}\\
\toprule
Lớp & Nằm ở & Ghi chú \\
\midrule
\endfirsthead
\toprule
Lớp & Nằm ở & Ghi chú \\
\midrule
\endhead
\bottomrule
\endfoot
Dịch vụ & \code{backend/app/api}, \code{backend/app/services}, \code{backend/app/schemas.py} & FastAPI tự kiểm tra dữ liệu vào theo kiểu đã khai báo và trả lỗi 422 khi sai, nên từng endpoint không phải tự kiểm tra từng trường \\
Cơ sở dữ liệu & \code{models.py} trong \code{backend/app/db}; các tệp di trú trong \code{backend/alembic} & Lược đồ chỉ đổi qua di trú Alembic; cột vector có chỉ mục HNSW \\
Suy luận mô hình & \code{face\_service.py}, \code{antispoof\_service.py} trong \code{backend/app/services} & Mô hình nạp ở lần dùng đầu tiên rồi giữ trong bộ nhớ cho các yêu cầu sau \\
Ứng dụng & \code{mobile/lib} & Cùng một mã nguồn build ra bản Web, Android và iOS \\
Triển khai & \code{docker-compose.yml}, \code{deploy/} & Một lệnh dựng cơ sở dữ liệu, máy chủ và giao diện \\
Kiểm thử & \code{backend/tests}, \code{mobile/test} & Phía máy chủ chạy trên PostgreSQL thật; phía ứng dụng dùng \code{flutter test} \\
\end{xltabular}}

\section{Các hệ thống liên quan}

Điểm danh tự động đã có nhiều lời giải thương mại và học thuật. Bảng
\ref{tab:sosanh} đối chiếu FAMS với ba nhóm giải pháp phổ biến.

{\small
\begin{xltabular}{\textwidth}{@{}>{\raggedright\arraybackslash}p{3.0cm} X X X@{}}
\caption{So sánh FAMS với các giải pháp điểm danh khác.}
\label{tab:sosanh}\\
\toprule
Tiêu chí & Điểm danh thủ công & Thẻ từ hoặc mã QR & FAMS \\
\midrule
\endfirsthead
\toprule
Tiêu chí & Điểm danh thủ công & Thẻ từ hoặc mã QR & FAMS \\
\midrule
\endhead
\bottomrule
\endfoot
Thời gian mỗi tiết & 5--10 phút & 1--2 phút & Dưới 1 phút, quét cả lớp \\
Chống điểm danh hộ & Không & Không, vì thẻ và mã chuyển tay được & Có, dựa
  trên đặc điểm sinh trắc \\
Bằng chứng kiểm chứng & Chữ ký giấy & Nhật ký quẹt thẻ & Ảnh bằng chứng kèm
  nhật ký sửa đổi \\
Chi phí thiết bị & Không & Đầu đọc cho mỗi phòng & Máy tính hoặc điện thoại
  sẵn có \\
Rủi ro về dữ liệu & Thấp & Trung bình & Cao, vì là dữ liệu sinh trắc, nên
  cần cơ chế đồng ý và phân quyền \\
\end{xltabular}}

{\small
\begin{xltabular}{\textwidth}{@{}>{\bfseries\raggedright\arraybackslash\leavevmode}p{2.6cm} >{\raggedright\arraybackslash\leavevmode}p{4.6cm} >{\raggedright\arraybackslash\leavevmode}X@{}}
\caption{Giải thích chi tiết Bảng~\ref{tab:sosanh}: ý nghĩa từng tiêu chí và căn cứ cho cột FAMS.}
\label{tab:sosanh-gt}\\
\toprule
Tiêu chí & Nghĩa là gì & Căn cứ cho cột FAMS \\
\midrule
\endfirsthead
\toprule
Tiêu chí & Nghĩa là gì & Căn cứ cho cột FAMS \\
\midrule
\endhead
\bottomrule
\endfoot
Thời gian mỗi tiết & Thời gian giảng viên dành cho việc điểm danh trong một tiết & Ước lượng từ thời gian xử lý mỗi khung hình ở Bảng~\ref{tab:hieunang}; chưa đo trên một lớp học thật \\
Chống điểm danh hộ & Có ngăn được một người điểm danh thay người khác hay không & Khuôn mặt không chuyển tay được; các lớp kiểm tra ở Bảng~\ref{tab:phongve} \\
Bằng chứng kiểm chứng & Căn cứ để xem lại khi có khiếu nại & Ảnh bằng chứng lưu \FngayLuuAnh{} ngày và nhật ký sửa điểm danh có lý do bắt buộc \\
Chi phí thiết bị & Thiết bị phải mua thêm cho mỗi phòng & Dùng webcam hoặc điện thoại sẵn có; máy chủ chỉ cần CPU \\
Rủi ro về dữ liệu & Mức nhạy cảm của dữ liệu thu thập & Dữ liệu sinh trắc, nên đồng ý là bắt buộc (Mục~\ref{sec:dongy}) và mọi truy cập đi qua phân quyền \\
\end{xltabular}}

Nhận xét rút ra là ưu thế của hướng nhận diện khuôn mặt chỉ giữ được nếu hệ
thống xử lý đúng ba việc: chống giả mạo, phân quyền truy cập ảnh, và cho
người dùng quyền kiểm soát dữ liệu sinh trắc của mình. Ba việc này định hình
phần lớn thiết kế được trình bày ở Chương~\ref{ch:thietke}.
